% R1 relation-map: structural diagram (no numeric axes), transcribed from
% create_relation_map() in figures/src/r1_figures.py -- same three rows, same
% base/target domain layout, same coordinates (scaled by a fixed unit length).
%
% ROW 1 REDRAWN (not transcribed): the script's row 1 was two blocks of prose
% stating the covering fact in words. Nothing was encoded on any visual
% channel, so the row conveyed strictly less than its own caption. It is
% replaced here by a complete worked instance at N=6, k=1, m=2 -- six dots,
% three drawn codewords, each shading the two placements it covers -- which
% is the entire counting argument of Theorem 1 at a size that fits in a
% corner. Numbers on it are closed-form: log2 C(6,1) - log2 C(2,1)
% = log2 6 - log2 2 = 1.585 bits, so 2 bits (4 codewords) suffice and 3 are
% used. Rows 2 and 3 are unchanged transcriptions.
%
% TAG CORRECTED (not silently): row 3 tagged the 0/16 falsification count
% [verified], contradicting both this fragment's own header tag and
% paper1.tex's three text sites, all of which say [internal]. The 0/16 count
% is a seeded Monte-Carlo sweep (16 budget points, 4000 placements each, an
% explicit 0.04 MC slack in the comparison), so it is [internal] by the
% paper's own stated definition. Corrected below.
% [internal, script r1_figures.py]
\begin{figure}[t]
\centering
\begin{tikzpicture}[x=0.78cm,y=0.78cm]
\node[align=center,font=\normalsize\bfseries,text=harborblue] at (6,12.35)
  {R1 --- the digest is a message; the floor is its length};

% Base domain (left): the combinatorial world
\draw[draw=harborblue,thick,fill=harborblue,fill opacity=0.12] (0.2,0.7) rectangle (3.8,10.6);
\node[align=center,font=\small\bfseries] at (2.0,11.4) {Base: where the $k$ criticals\\could be hiding};

% Target domain (right): the digest channel
\draw[draw=seagreen,thick,fill=seagreen,fill opacity=0.12] (8.2,0.7) rectangle (11.8,10.6);
\node[align=center,font=\small\bfseries] at (10.0,11.4) {Target: the digest that must\\name where to look};

% Row 1 --- drawn worked instance at N=6, k=1, m=2 (see header note).
\node[align=center,font=\scriptsize,text width=2.5cm] at (2.0,9.60) {$N{=}6$ artifacts;\\$k{=}1$ is critical};
\foreach \i/\x in {1/0.75,2/1.25,3/1.75,4/2.25,5/2.75,6/3.25}{
  \fill[harborblue] (\x,8.55) circle (0.09);
  \node[font=\tiny,text=black!70] at (\x,8.15) {\i};
}
\node[align=center,font=\scriptsize,text width=2.5cm] at (2.0,7.45) {$\binom{6}{1}=6$ placements};
% Target: the codebook, three m=2 codewords, each shading what it covers.
\node[align=center,font=\scriptsize,text width=2.6cm] at (10.0,9.60) {$2$ bits $\Rightarrow$ $4$ codewords,\\each an $m{=}2$ subset};
\foreach \y/\c in {8.80/00, 8.30/01, 7.80/10}{
  \node[font=\tiny\ttfamily,text=black!70] at (8.42,\y) {\c};
  \foreach \x in {8.85,9.35,9.85,10.35,10.85,11.35}{
    \fill[black!25] (\x,\y) circle (0.075);
  }
}
\draw[seagreen,thick,rounded corners=2pt,fill=seagreen,fill opacity=0.28] (8.70,8.66) rectangle (9.50,8.94);
\draw[seagreen,thick,rounded corners=2pt,fill=seagreen,fill opacity=0.28] (9.70,8.16) rectangle (10.50,8.44);
\draw[seagreen,thick,rounded corners=2pt,fill=seagreen,fill opacity=0.28] (10.70,7.66) rectangle (11.50,7.94);
\node[align=center,font=\scriptsize,text width=2.6cm] at (10.0,6.85) {each covers $\binom{2}{1}{=}2$ of the $6$; $3$ tile all $6$, the $4$th is slack};
\draw[<->,>=Stealth,thick,shipred] (3.9,7.93) -- (8.1,7.93);
\node[align=center,font=\small\bfseries,text=shipred] at (6.0,9.75) {one digest selects\\ONE covering};

% Row 2
\node[align=center,font=\scriptsize] at (2.0,5.00) {A digest of $B$ literal bits\\is a message with\\$2^B$ distinguishable\\values --- no more};
\node[align=center,font=\scriptsize] at (10.0,5.00) {Decoder: shared table\\of $m$-subsets; each\\message names one\\$m$-subset to open};
\draw[<->,>=Stealth,thick,shipred] (3.9,4.73) -- (8.1,4.73);
\node[align=center,font=\small\bfseries,text=shipred] at (6.0,6.55) {zero-miss needs\\$B \ge \log_2\binom{N}{k} - \log_2\binom{m}{k}$};

% Row 3
\node[align=center,font=\scriptsize] at (2.0,1.80) {Boundary instance:\\$N{=}60$, $k{=}2$ --- oracle /\\noisy / random encoders\\run against the bound};
\node[align=center,font=\scriptsize] at (10.0,1.80) {$m{=}8$ opened $\Rightarrow$\\$B^\star = 5.98$ bits [verified];\\$0/16$ floor violations\\{[internal, \texttt{a7\_experiment.py}]}};
\draw[<->,>=Stealth,thick,shipred] (3.9,1.53) -- (8.1,1.53);
\node[align=center,font=\small\bfseries,text=shipred] at (6.0,3.35) {the floor holds\\exactly};

\node[align=center,font=\scriptsize,text width=10cm] at (6.0,0.25)
  {\itshape The map carries a combinatorial fact, not scenery: each flagged $m$-set can only ever cover $\binom{m}{k}$ of the $\binom{N}{k}$ ways $k$ criticals could sit among $N$. Row 1 is the whole theorem in one instance --- $\log_2 6-\log_2 2=1.585$ bits, so $2$ bits and $3$ of the $4$ available codewords cover every placement [verified].};
\end{tikzpicture}
\caption{The relation-map for Theorem 1. A digest is a $B$-bit message selecting one of the decoder's covering $m$-subsets (base domain: a shared codebook); each covering accounts for $\binom{m}{k}$ of $\binom{N}{k}$ placements, which is the whole proof. The map is of \emph{relations}, not surface features: from it the reader can re-derive the floor without the text.}
\label{fig:r1rel}
\end{figure}
